% Ingested from kunci_jawaban (tugas_topik_3.tex)
% shared
\begin{aligned}
x \to -1^{-}: \quad f(x) = x + 1 \to 0 \\
x = -1: \quad f(-1) = (-1)^2 = 1 \\
x \to 2^{-}: \quad f(x) = x^2 \to 4 \\
x = 2: \quad f(2) = 2 - 2 = 0 \\
D_f = (-\infty, -1) \cup [-1, 2) \cup [2, \infty) = \mathbb{R} \\
x < -1 \;\Rightarrow\; x + 1 < 0 \\
-1 \le x < 2 \;\Rightarrow\; 0 \le x^2 < 4 \\
x \ge 2 \;\Rightarrow\; 2 - x \le 0
\end{aligned}

% role: concept
\begin{aligned}
x \to -1^{-}: \quad f(x) = x + 1 \to 0 \\
x = -1: \quad f(-1) = (-1)^2 = 1 \\
x \to 2^{-}: \quad f(x) = x^2 \to 4 \\
x = 2: \quad f(2) = 2 - 2 = 0 \\
Grafik di x = -1 (dari ketinggian 0 ke 1 ) dan di x = 2 (dari ketinggian 4 ke 0 ). \\
grafik melompat di titik ujung selang (titik kosong vs penuh)
\end{aligned}

% role: value_table
\begin{aligned}
(-3, -2) \\
(-2, -1) \\
(-1, 1) \\
(0, 0) \\
(1, 1) \\
(1{,}5, 2{,}25) \\
(2, 0) \\
(3, -1) \\
(4, -2)
\end{aligned}

% role: domain
\begin{aligned}
D_f = (-\infty, -1) \cup [-1, 2) \cup [2, \infty) = \mathbb{R} \\
x < -1 \;\Rightarrow\; x + 1 < 0 \\
-1 \le x < 2 \;\Rightarrow\; 0 \le x^2 < 4 \\
x \ge 2 \;\Rightarrow\; 2 - x \le 0
\end{aligned}

% role: range
\begin{aligned}
R_f = (-\infty, 0) \cup [0, 4) \cup (-\infty, 0] = (-\infty, 4).
\end{aligned}

% role: figure

% final answer
D_f = \mathbb{R} ; R_f = (-\infty, 4)
